% Tento soubor obsahuje definici prostředí pro popis případů užití.
% 
% Použití:
% \begin{UseCaseTable}                                        %Začátek tabulky případů užití
%   {Přihlášení uživatele}                                    %Název případu užití
%   {UC-login}                                                %Identifikace případu užití
%   {FR-01}                                                   %Identifikace funkčního požadavku
%   {Uživatel se přihlásí do systému}                         %Cíl případu užití
%   {Uživatel}                                                %Primární aktér(ři)
%   {Autentizační služba}                                     %Pomocný aktér(ři)
%   {Uživatel není přihlášen}                                 %Vstupní podmínky
%   {Uživatel je přihlášen}                                   %Výstupní podmínky
%
% \UseCaseScenario{Základní scénář}                           %Vymezení zda li jde o základní nebo alternativní scénář (dle MMSP metodiky)
% \UseCaseStep{uživatel}{zadá přihlašovací údaje}             %Role (Systém/aktér) a akce (co dělá)
%
% \UseCaseScenario{Alternativní scénář}
% \UseCaseStepCustom{A1}{systém}{neplatné přihlašovací údaje} %Custom step pro alternativní scénář, kde krok není číslován explicitně.
%
% \end{UseCaseTable}                                          %Konec tabulky případů užití

\newcounter{ucstep}

\newenvironment{UseCaseTable}[8]
{
\footnotesize
\setlength{\tabcolsep}{3pt}
\renewcommand{\arraystretch}{1.1}
\setcounter{ucstep}{0}

\begin{longtable}{p{4cm} p{2cm} p{3cm} p{5cm}}
\caption{Případ užití: #1}\label{tab:usecase:#2}\label{#2}\tabularnewline

\toprule
\textbf{Název případu užití} & \multicolumn{3}{l}{#1} \\
\midrule

\textbf{Identifikace případu užití} & \multicolumn{3}{p{10cm}}{#2} \\
\textbf{Identifikace funkčního požadavku} & \multicolumn{3}{p{10cm}}{\hyperref[#3]{#3}} \\
\textbf{Cíl případu užití} & \multicolumn{3}{p{10cm}}{#4} \\
\textbf{Primární aktér(ři)} & \multicolumn{3}{p{10cm}}{#5} \\
\textbf{Pomocný aktér(ři)} & \multicolumn{3}{p{10cm}}{#6} \\
\textbf{Vstupní podmínky} & \multicolumn{3}{p{10cm}}{#7} \\
\textbf{Výstupní podmínky} & \multicolumn{3}{p{10cm}}{#8} \\
% \textbf{Základní scénář} & \textbf{Krok} & \textbf{Role} & \textbf{Akce} \\
}
{
\bottomrule
\end{longtable}
}

\newcommand{\UseCaseStep}[2]{%
\stepcounter{ucstep}
& \arabic{ucstep} & #1 & #2 \\
}

\newcommand{\UseCaseScenario}[1]{%
\midrule
\textbf{#1} & \textbf{Krok} & \textbf{Role} & \textbf{Akce} \\
\setcounter{ucstep}{0}
}

\newcommand{\UseCaseStepCustom}[3]{%
& #1 & #2 & #3 \\
}
